\documentclass[12pt]{article}
\usepackage{amsmath,amssymb}

\begin{document}
\section*{{2025 AMC 12B Problems/Problem 13}}
Source: AoPS Wiki

\subsection*{Solution}
Create an arbitrary six slice circular diagram. The first slice has 3 options, and the second has 2, third has 2, fourth has 2, fifth has 2, and the sixth has only 1 color to be chosen. Understand that if \( \mathbb{X} \) is the set containing all possible cases of the circle, then \( \mathbb{X} \) has a correlation to another set, call it \( \mathbb{K} \), which contains all the colors. This can be quite easily sought as a rotation of one element in \( \mathbb{X} \) that accurately maps to another element in \( \mathbb{X} \) is considered the same. In simpler terms, to account for overlapping cases in \( \mathbb{X} \), one must divide by 2 (as two colorings that are rotated are considered different, and we don't want that as it defeats the purpose of "unique"). Then, multiplying out we have \( 3 \times 2^4 = 48 \). Dividing by 2 to account for overlap in \( \mathbb{X} \), we get \( 48/2 \) or $\boxed{\textbf{(D) } 24}$ ~Pinotation

\end{document}
